
% Thesis Abstract -----------------------------------------------------


%\begin{abstractslong}    %uncommenting this line, gives a different abstract heading
\begin{abstracts}        %this creates the heading for the abstract page

El presente trabajo de tesis se enfoca en el diseño y control de la aeronave de ala-integrada EHEKATL, conocida como EHEKATL BWB (\textit{Blended Wing Body}). Los principales capítulos abordan su diseño, modelado dinámico, diseño de controladores y las simulaciones de prueba realizadas en el entorno del simulador de vuelo X-PLANE.

El proceso de diseño de la aeronave EHEKATL BWB se inicia utilizando el software de código abierto XFLR5 para obtener los parámetros aerodinámicos cruciales para la estabilidad y control de la aeronave. Posteriormente, se procede a la configuración detallada en las extensiones especializadas de X-PLANE: AIRFOIL MAKER se emplea en el diseño de los perfiles aerodinámicos seleccionados, mientras en PLANE MAKER se diseña el modelo completo de la aeronave considerando su geometría, aerodinámica, peso y equilibrio, sistemas y características de rendimiento.

Una vez definido el diseño, se desarrolla el modelo dinámico no lineal basado en la mecánica Newtoniana.

Se presenta el diseño de los controladores Backstepping y Backstepping Adaptable. Estas técnicas de control no solo permiten abordar la complejidad inherente de la dinámica no lineal de la aeronave, sino que también ofrecen una mayor robustez y adaptabilidad en situaciones donde los parámetros del sistema pueden ser desconocidos o variables. Esta capacidad de adaptación es de suma importancia en entornos dinámicos y cambiantes como los de vuelo. Se presentan una serie de simulaciones numéricas para comprobar el desempeño de los controladores propuestos. 

Adicionalmente, se presentan una serie de simulaciones cuasi-reales en el entorno de X-PLANE. 
Los resultados obtenidos en las simulaciones de vuelo con X-PLANE validan el desempeño de la aeronave EHEKATL BWB con la aplicación del controlador Backstepping. Los resultados muestran un desempeño eficiente en el controlador Backstepping tanto en la dinámica longitudinal como en la lateral-direccional.

\end{abstracts}
%\end{abstractlongs}
% ----------------------------------------------------------------------